\documentclass[10pt,a4paper]{amsart}
\usepackage[margin=19mm]{geometry}
\usepackage{amsmath,amssymb,mathtools}
\usepackage[scr=rsfs,cal=euler]{mathalfa}
\usepackage{xfrac}
\usepackage[unicode,hidelinks,bookmarks=false]{hyperref}
\newcommand{\ie}{i.e.\ }
\newcommand{\cf}{cf.\ }
\newcommand{\resp}{respectively\ }
\newcommand{\Subsection}{\subsection}
\providecommand{\nobreakdash}{\nobreak\hbox{-}}
\setlength{\emergencystretch}{2em}
\multlinegap0pt
\usepackage{amssymb}
\usepackage{enumerate}
\newtheorem{lemma}{Lemma}
\newcommand{\conservativepart}{\mathfrak{Cons}}
\newcommand{\convgenus}{{\boldsymbol \xi}}
\newcommand{\dissipativepart}{\mathfrak{Diss}}
\newcommand{\ext}{{\rm Ext}}
\newcommand{\teich}{\mathcal{T}}
\title{Hopf decomposition of the actions of subgroups of the mapping class group}
\author{Hideki Miyachi}
\date{Source: arXiv:2609.10175v1 (CC BY 4.0)}
\begin{document}
\maketitle

\noindent\emph{Source and licence.} Hopf decomposition of the actions of subgroups of the mapping class group. Version arXiv:2609.10175v1 (CC BY 4.0), distributed by arXiv under the Creative Commons Attribution 4.0 International licence, http://creativecommons.org/licenses/by/4.0/, as recorded in the arXiv licence metadata for this exact version and shown on the abstract page https://arxiv.org/abs/2609.10175v1. Full licence notice: \texttt{licenses/arxiv-2609.10175-CC-BY-4.0.txt}.\par\smallskip

\noindent\textit{Editorial scope.} Counted unit: the article's lemma lem:Cons-Diss-ElementaryProperty (source0.tex lines 2166--2176). The article's complete original proof of it is delivered with it (source0.tex lines 2178--2226, one \texttt{proof} environment of 49 lines). No mathematics is written, repaired or re-derived: the statement and the proof are the article's own lines, in the article's own notation, and no retained line is substituted or re-flowed.  Retained uncounted as the source-local context the proof needs: lines 659--662.  Nothing was omitted from any retained span, so the package carries no omission record.  No retained line contains a citation, so the wrapper carries no bibliography and no bibliography entry.  No retained line contains a figure environment, an image-inclusion command or drawing code, so no image is delivered and the package declares no asset dependency.  Editorial formatting only, in addition: an AMS \texttt{amsart} preamble that restates the article's own declarations and macros used by the retained lines, namely the environments \texttt{lemma} on separate counters, together with the standard packages those lines need.  The retained environments print in this document as Lemma 1 in order of appearance, whereas the article prints the same items with the numbering of its own full document.  The complete native source of the article as distributed in the arXiv e-print for this exact version is bundled unchanged as \texttt{sources/209881/source0.tex}.
\begin{equation}
\label{eq:QC-inv_extremal_length}
e^{-2d_T(x,y)}\le \frac{\ext_x(\lambda)}{\ext_y(\lambda)}\le e^{2d_T(x,y)}
\end{equation}

\begin{lemma}
\label{lem:Cons-Diss-ElementaryProperty}
Let $G$ be a subgroup of the Teichm\"uller modular group.
\begin{itemize}
\item[(a)]
The sets $\conservativepart(G)$ and $\dissipativepart(G)$ are independent of the choice of basepoint $x\in \teich_{g,m}$.
\item[(b)]
$\conservativepart(G)$ and $\dissipativepart(G)$ are invariant under
the action of $G$.
\end{itemize}
\end{lemma}

\begin{proof}
\textbf{(a).}
By \eqref{eq:QC-inv_extremal_length} and the invariance of $d_T$ under the action of $G$, we have
\[
e^{-4\convgenus d_T(x,y)}
\left(
\frac{\ext_x(\lambda)}
{\ext_{[\omega](x)}(\lambda)}
\right)^{\convgenus}
\le
\left(
\frac{\ext_y(\lambda)}
{\ext_{[\omega](y)}(\lambda)}
\right)^{\convgenus}
\le
e^{4\convgenus d_T(x,y)}
\left(
\frac{\ext_x(\lambda)}
{\ext_{[\omega](x)}(\lambda)}
\right)^{\convgenus}
\]
for all $x,y\in\teich_{g,m}$ and $[\omega]\in G$.
Summing over $[\omega]\in G$ shows that convergence or divergence
of the defining series is independent of the choice of basepoint.

\medskip
\noindent
\textbf{(b).}
For $[\eta]\in G$, the equivariance of extremal length gives
\begin{align*}
&\sum_{[\omega]\in G}
\left(
\frac{\ext_x([\eta](\lambda))}
     {\ext_{[\omega](x)}([\eta](\lambda))}
\right)^{\convgenus} =
\sum_{[\omega]\in G}
\left(
\frac{\ext_{[\eta]^{-1}(x)}(\lambda)}
     {\ext_{[\eta]^{-1}\circ[\omega](x)}(\lambda)}
\right)^{\convgenus}=
\sum_{[\gamma]\in G}
\left(
\frac{\ext_{[\eta]^{-1}(x)}(\lambda)}
     {\ext_{[\gamma]\circ[\eta]^{-1}(x)}(\lambda)}
\right)^{\convgenus},
\end{align*}
where $[\gamma]=[\eta]^{-1}\circ[\omega]\circ[\eta]$.
Since the convergence or divergence of the defining series is independent of the choice of basepoint, this proves~(b).
\end{proof}

\end{document}
